\documentclass[leqno, a4paper, 12pt]{jsarticle}
\usepackage{amssymb}
\usepackage{amsthm}
\usepackage{amsmath}
\usepackage{comment}
	%可換図式用
		%\usepackage[all]{xy}
	%重くなるので使わいときはコメント化しておく。
%定理環境 thmはあまり使わずpropをなるべく使う。
%主張は基本的に証明の中で使い、そのため番号を付けない。
%注意環境を付けてみた。
\newtheorem{thm}{定理}
\newtheorem{prop}[thm]{命題}
\newtheorem{cor}[thm]{系}
\newtheorem{lem}[thm]{補題}
\newtheorem{df}[thm]{定義}
\newtheorem{exam}[thm]{例}
\newtheorem{fact}[thm]{事実}
\newtheorem*{claim}{主張}
\newtheorem*{rmk}{注意}
\renewcommand{\proofname}{{\bf 証明}}
%矢印たち。
%\toは\rightarrowと同じ。\getsは\leftarrowと同じ。同値は\iff
%上向き矢はuparrow逆はdownarrow
\newcommand{\To}{\Longrightarrow}
\newcommand{\Gets}{\Longleftarrow}
%黒板太字
\newcommand{\nn}{\mathbb{N}}
\newcommand{\zz}{\mathbb{Z}}
\newcommand{\qq}{\mathbb{Q}}
\newcommand{\rr}{\mathbb{R}}
\newcommand{\cc}{\mathbb{C}}
%無限大
\newcommand{\mgn}{\infty}
%ギリシャン
\newcommand{\dpst}{\displaystyle}
\newcommand{\ep}{\varepsilon}
\newcommand{\al}{\alpha}
\newcommand{\be}{\beta}
\newcommand{\de}{\delta}
\newcommand{\la}{\lambda}
\newcommand{\La}{\Lambda}
\newcommand{\ga}{\gamma}
\newcommand{\lil}{\la\in \La}
%商集合
\newcommand{\qt}[2]{#1/\! #2}
%enumerate環境
\renewcommand{\labelenumi}{{\rm (\arabic{enumi})}}
%以降alignじゃなくてenumerate を使え。
%なんか演算
\DeclareMathOperator{\card}{card}
\DeclareMathOperator{\supp}{supp}
%なんか演算２
\DeclareMathOperator{\bd}{Bdry}
\DeclareMathOperator{\cl}{CL}
\DeclareMathOperator{\nai}{Int}
\DeclareMathOperator{\di}{diam}

%差集合は\setminusで統一する。等号の否定は\neq
%真部分集合は\subsetneqq　偏微分は\partial
%ディスプレイスタイル。\dfracでディスプレイ分数
%空集合は\Oを使う！！！！！！！！！！！！

\title{Radon-Nikodymの定理}
\author{ナンブキトラ}
\date{\today}
			\begin{document}
\maketitle

この文章ではRadon-Nikodymの定理を証明する．

\begin{df}
$(X, \mathfrak{M})$
を測度空間とし，$\mu$
をその上の測度とする．
このとき
$L^{p}(\mu)$
で$\mu$に関する$L^{p}$関数全体を表す．
つまり$f\in L^{p}(\mu)$であるとは
$f: X\to \rr$であって可測であり，
\[
\int_{X} |f|^{p}\ d\mu<\infty
\]
を満たすもの全体である．

空間$L^{2}(\mu)$はヒルベルト空間になることに注意しよう．
\end{df}


\begin{lem}\label{lem:lemlem}
$(X, \mathfrak{M})$
を測度空間とし，$\alpha$
をその上の$\sigma$有限な測度とする．
このとき可測関数$u: X\to (0, 1)$で
$y\in L^{1}(\alpha)$
を満たすものが存在する．
\end{lem}
\begin{proof}

$\{A_{i}\}_{i\in \zz_{\ge 0}}$
を
\begin{enumerate}
\item 各$A_{i}$は可測
\item 
$\bigcup_{i\in \zz_{\ge 0}}A_{i}=X$
\item $i\neq j$ならば$A_{i}\cap A_{j}=\emptyset$
\item 
と任意の$i$について$\alpha(A_{i})<\infty$
\end{enumerate}
を満たすものとする．
そして$u: X\to (0, 1)$を
\[
u(x)=\sum_{n=1}^{\infty}\frac{\chi_{A_{n}}(x)}{2^{n}(1+\alpha(A_{n}))}
\]
と定義すれば求める条件は満たされる．
\end{proof}

\begin{df}
$(X, \mathfrak{M})$
を可測空間とし，$\alpha$, $\beta$
をその上の測度とする．
このとき$\alpha$が$\beta$に関して絶対連続であるとは，
$\beta(A)=0$となる$A\in \mathfrak{M}$
について$\alpha(A)=0$となることである．このとき
$\alpha\ll \beta$と表し，
$\alpha\ll\beta$かつ$\beta\ll\alpha$になるとき$\alpha \equiv \beta$と表す．この$\equiv$は同値関係になる．
また，ある$A\in \mathfrak{M}$が存在して
$\alpha(A)=0$と$\beta(X\setminus A)=0$
を満たすときに$\alpha$は$\beta$について特異的であるといい，$\alpha \perp \beta$と書く．
この$\perp$は同値関係になる．
\end{df}

\begin{lem}
$(X, \mathfrak{M})$
を測度空間とし，$\alpha, \beta, \gamma$
をその上の測度とする．
もしも$\alpha\equiv \beta$
が成り立つならば以下が成り立つ．
\begin{enumerate}
\item 
$\beta\ll \gamma$ならば
$\alpha\ll \gamma$が成り立つ. 
\item 
$\gamma\ll \beta$
ならば
$\gamma\ll \alpha$が成り立つ. 
\item 
$\beta\perp \gamma$
ならば
$\alpha\perp \gamma$
が成り立つ．
\end{enumerate}
\end{lem}
\begin{proof}
証明は省略する．
\end{proof}


\begin{lem}\label{lem:hyougen}
$(X, \mathfrak{M})$
を測度空間とし，$\alpha$
をその上の測度とし，
$\phi\in L^{1}(\alpha)$とする．
そして測度$\beta$を
\begin{equation}\label{eq:1}
\beta(E)=\int_{E}\phi\ d\alpha
\end{equation}
と定義すると，
$\beta\ll \alpha$であり，
任意の$f\in L^{1}(\beta)$について
\[
\int_{X}f\ d\beta
=
\int_{X}f\phi\ d\alpha
\]
が成り立つ．
また測度$\beta$についてある$\phi$が存在して(\ref{eq:1})
のように表されるのならば$\beta\ll \alpha$となる．
\end{lem}
\begin{proof}
$\beta$の定義から任意の$E\in \mathfrak{M}$について
\[
\int_{X}\chi_{E}\ d\beta=\int_{X} \chi_{E}\phi\ d\alpha
\]
が成り立つ．
つまり特性関数については補題の結論は成り立つ．
このことから単関数についても成り立つし，
単関数近似から一般の可測関数についても成り立つ．
後半は明らかである．
\end{proof}

\begin{lem}[リースの表現定理]\label{lem:riesz}
$(H, \langle \cdot, \cdot \rangle_{H})$をヒルベルト空間とする．
このとき任意の連続な線形写像
$f: H\to \rr$に対して$a\in H$ご存在して
任意の$x\in H$について
$f(x)=\langle x, a\rangle_{H}$が成り立つ．
\end{lem}


以下の定理がラドン・ニコディムの定理である．
\begin{thm}
$(X, \mathfrak{M})$
を測度空間とし，$\mu$，$\nu$
をその上の$\sigma$有限な測度とする．
このとき$(X, \mathfrak{M})$
上の測度$\nu_{a}$と
$\nu_{s}$が存在して以下を満たす．
\begin{enumerate}
\item 
$\nu_{s}\perp \mu$
と$\nu_{a}\ll \mu$を満たす．
\item 
$(1)$のような分解は一意的である．
\item 
可測関数
$D: X\to \rr$
が存在して
任意の$A\in \mathfrak{M}$について
$\nu_{a}(A)=\int_{A}D\ d\mu$
が成り立つ．
\end{enumerate}
\end{thm}
\begin{proof}
$(2)$はすぐにわかる．

$(1)$を示そう．
補題 \ref{lem:lemlem}から
$s, t: X\to (0, 1)$を
$s\in L^{1}(\mu)$
で$t\in L^{1}(\nu)$
を満たすものとする．
そして$(X, \mathfrak{M})$上の
測度$\lambda$, 
$\theta$を
\[
\lambda(A)=
\int_{A} s\ d\mu
\]
\[
\theta(A)=
\int_{A} t\ d\nu
\]
と定義する．
$\lambda$, 
$\theta$
は有限測度であることに注意しよう．
また
\begin{align}
\mu(A)=
\int_{A} \frac{1}{s}\ d\lambda\\
\nu(A)=
\int_{A} \frac{1}{t}\ d\theta\label{al:4}
\end{align}
なので
$\lambda\equiv \mu$と$\theta\equiv \nu$に注意しよう．
$\phi=\lambda+\theta$と定義する．
すると，コーシーシュワルツの不等式から
$f\in L^{2}(\phi)$
に対して
\[
\left|
\int_{X}f\ d\theta
\right|
\le 
\int_{X}|f|\ d\theta
\le 
\int_{X}|f|\ d\phi
\le 
\phi(X)\cdot 
\sqrt{
\int_{X}|f|^{2}\ d\phi}
<\infty
\]
を得る．ここで$s\in L^{1}(\mu)$, $t\in L^{1}(\nu)$なので
$\phi(X)=\lambda(X)+\theta(X)<\infty$
であることを用いている．
よって
$f\in L^{2}(\phi)$
に対して
$\int_{X}f\ d\theta$
と定義される線形形式は
$L^{2}(\phi)$上の有界作用素となる．
よって
$L^{2}(\phi)$
がヒルベルト空間であることから
補題 \ref{lem:riesz}を使うと
任意の$f\in L^{2}(\phi)$
に対して
\begin{equation}\label{eq:hyo}
\int_{X}f\ d\theta
=
\int_{X}f g \ d\phi
\end{equation}
となる$g\in L^{2}(\phi)$
が存在する．
さて$\phi(E)>0$となる可測集合
$E$について，$f=\chi_{E}$とすると
\[
\int_{E}g\ d\phi=\theta(E)
\]
となる．
よって特に
\begin{equation}\label{eq:hutou}
0\le \frac{1}{\phi(E)}\int_{E}g\ d\phi\le 1
\end{equation}
となる．
ここで，背理法のために$S=\{\, x\in X\mid 1<g(x)\, \}$
が$\phi$に関して正の測度を持つと仮定しよう．
$\epsilon$に対して
$T_{\epsilon}=\{\, x\in X\mid 1+\epsilon<g(x)\, \}$
と仮定すると
\[
S=\bigcup_{\epsilon \in \qq_{>0}}T_{\epsilon}
\]
が成り立つ．よって
ある$\epsilon\in \qq_{>0}$
が存在し，
$T_{\epsilon}=\{\, x\in X\mid 1+\epsilon<g(x)\, \}$
が$\phi$に関して正の測度をもつ．
よって式 (\ref{eq:hutou})から
\[
1<1+\epsilon\le \frac{1}{\phi(T_{\epsilon})}\int_{T_{\epsilon}}g\ d\phi\le 1
\]
となり矛盾する．
よってほとんど至る所$g\le 1$である．
同様にほとんど至る所$0\le g$であることがわかる．
つまりほとんど至る所$0\le g(x)\le 1$である．
さて式 (\ref{eq:hyo})を変形して
\begin{equation}\label{eq:hyo2}
\int_{X}f(1-g)\ d\theta
=
\int_{X}f g \ d\lambda
\end{equation}
を得る．
ここで
$A=\{\, 
x\in X\mid 
0\le g(x)<1
\, \}$, 
$S=
\{\, 
x\in X\mid 
g(x)=1
\, \}
$
と定義して
$\theta_{a}(E)=\theta(E\cap A)$, 
$\theta_{s}(E)=\theta(E\cap S)$
と定義する．
式 (\ref{eq:hyo2})において
$f=\chi_{S}$とすれば，$\lambda(S)=0$を得るので
$\theta_{s}\perp \lambda$
がわかる．
さて，
$h$を
\[
h(x)=
\begin{cases}
\frac{g(x)}{1-g(x)} & \text{if $x\in A$}\\
0 & \text{その他 }
\end{cases}
\]
と定義する．
各$n\in \zz_{\ge 0}$, $E\in \mathfrak{M}$について
式 (\ref{eq:hyo2})において
$f=(1+g+\dots+g^{n})\chi_{E\cap A}$
とおくと
\[
\int_{E\cap A}1-g^{n+1}\ d\theta
=
\int_{E\cap A}g(1+g+\dots+g^{n})\ d\lambda
\]
を得る．この式で
$n\to \infty$とすれば
\[
\theta_{a}(E)=\theta(E\cap A)=\int_{E}h \ d\lambda
\]
を得る．よって$\theta_{a}\ll \lambda$である．
ここで
\[
\nu_{a}(A)=\int_{A}\frac{1}{t}\ d\theta_{a}
\]
\[
\nu_{s}(A)=\int_{A}\frac{1}{t}\ d\theta_{s}
\]
とすれば$\nu_{a}\equiv \theta_{a}$と
$\nu_{s}\equiv \theta_{s}$であり
式 (\ref{al:4})から$\nu_{a}+\nu_{s}=\nu$
となり，$\lambda\equiv \mu$
から$\nu_{a}\ll \mu$と$\nu_{s}\perp \mu$
が成り立つ．




$(3)$を示そう．
今
\[
\nu(E)
=
\int_{E}\frac{1}{t}\ d\theta
\]
なので
\[
\nu_{a}(E)
=
\int_{E}\frac{1}{t}\ d\theta_{a}
=\int_{E}\frac{h}{t}\ d\lambda
=\int_{E}\frac{sh}{t}\ d\mu
\]
となるから
\[
D=
\frac{sh}{t}
\]
とすれば良い．
\end{proof}


























	
\begin{thebibliography}{9}
\bibitem{1}
W. Rudin, 
\textit{Real and complex analysis}, 
3rd ed.,  McGRAW--Hill. 
\end{thebibliography}




			\end{document}



\begin{comment}
[集合のやつは$\mid$を使う。]
[真なる包含関係]
\subsetneqq
[可換図式]
\[\xymatrix{
&   & (2) \ar@{=>}[rd]&  &  &  & (5) \ar@{=>}[rd]&\\ 
& (1)\ar@{=>}[ru] \ar@{=>}[rd]&  &(4)\ar@{=>}[r]&(7)& (1)\ar@{=>}[ru] \ar@{=>}[rd]& & (7)&\\ 
&   & (3) \ar@{=>}[ur]&  &  &  &(6) \ar@{=>}[ur]&
}\]

[\bigin \endを使わない方法]

{\prop[aa] aaa}
で
\begin{prop}[aa]
aaa
\end{prop}
と同じ\df \thm \proof でも同様

[直和]
$\amalg$

[同値関係でよく使う波線]
\sim

[引数付きの新命令]
\newcommand{\prn}[1]{\|#1\|_p}

[番号のリセット]

\setcounter{equation}{0}


[字体の変更]

{\rm hogehoge}と使う。

\rm ローマン
\it イタリック
\sl 斜体

[supp]

\newcommand{\supp}{\text{{\rm supp}}}

[中央揃えの改行]

	\begin{eqnarray*}
		hoge1&=&hogehoge
		hoge2&=&hogehofe

	\end{eqnarray*}

&&の部分で揃えられる。


[\tag]
\begin{equation}
f(x) = ax^2 + bx + c \tag{1.2.3}
\end{equation}



[場合分け]

	\begin{cases}
		hoge1 & \text{状況１}\\
		hoge2 & \text{状況２}\\
		・
		・
		・
	\end{cases}



\end{comment}





